% ============================================================================
% (c) 2025 Mohamed EL AFRIT -- IPSSI
% Licence CC BY-NC-SA 4.0
% https://creativecommons.org/licenses/by-nc-sa/4.0/deed.fr
%
% FICHE DE CORRECTIONS DETAILLEES -- Jour 1
% Donnees et representation mathematique
% ============================================================================
% ============================================================================
% © 2025 Mohamed EL AFRIT — IPSSI
% Ce contenu est distribué sous licence Creative Commons BY-NC-SA 4.0
% https://creativecommons.org/licenses/by-nc-sa/4.0/deed.fr
% ============================================================================
%
% TEMPLATE COMMUN — Mathématiques appliquées à l'Intelligence Artificielle
% Ce fichier est le préambule partagé de TOUS les documents LaTeX du projet.
%
% Utilisation :
%   % ============================================================================
% © 2025 Mohamed EL AFRIT — IPSSI
% Ce contenu est distribué sous licence Creative Commons BY-NC-SA 4.0
% https://creativecommons.org/licenses/by-nc-sa/4.0/deed.fr
% ============================================================================
%
% TEMPLATE COMMUN — Mathématiques appliquées à l'Intelligence Artificielle
% Ce fichier est le préambule partagé de TOUS les documents LaTeX du projet.
%
% Utilisation :
%   % ============================================================================
% © 2025 Mohamed EL AFRIT — IPSSI
% Ce contenu est distribué sous licence Creative Commons BY-NC-SA 4.0
% https://creativecommons.org/licenses/by-nc-sa/4.0/deed.fr
% ============================================================================
%
% TEMPLATE COMMUN — Mathématiques appliquées à l'Intelligence Artificielle
% Ce fichier est le préambule partagé de TOUS les documents LaTeX du projet.
%
% Utilisation :
%   \input{template_ipssi}          % Importer ce préambule
%   \jourNumero{1}                  % Définir le numéro du jour (1, 2, 3 ou 4)
%   \jourTheme{Données et représentation mathématique}
%   \begin{document}
%   \pagegarde{Fiche de synthèse}{1}{Données et représentation mathématique}
%   ...
%   \end{document}
% ============================================================================

% --- Classe de document ---
\documentclass[11pt,a4paper]{article}

% ============================================================================
% ENCODAGE ET LANGUE
% ============================================================================
\usepackage[utf8]{inputenc}
\usepackage[T1]{fontenc}
\usepackage[french]{babel}

% ============================================================================
% MATHÉMATIQUES
% ============================================================================
\usepackage{amsmath,amssymb,amsfonts,mathtools}

% Commandes mathématiques personnalisées (conventions du cours)
\newcommand{\vx}{\mathbf{x}}          % vecteur x
\newcommand{\vw}{\mathbf{w}}          % vecteur w (poids)
\newcommand{\vy}{\mathbf{y}}          % vecteur y
\newcommand{\mX}{\mathbf{X}}          % matrice X
\newcommand{\mW}{\mathbf{W}}          % matrice W
\newcommand{\yhat}{\hat{y}}           % prédiction
\newcommand{\pscal}[2]{\langle #1, #2 \rangle} % produit scalaire
\newcommand{\gradJ}{\nabla J(\boldsymbol{\theta})} % gradient de J

% ============================================================================
% MISE EN PAGE
% ============================================================================
\usepackage[margin=2.5cm]{geometry}
\usepackage{setspace}
\setstretch{1.15}                      % interligne légèrement aéré

% ============================================================================
% COULEURS
% ============================================================================
\usepackage[dvipsnames,svgnames,x11names]{xcolor}

% Palette principale du projet
\definecolor{ipssiBleu}{HTML}{3B82F6}       % Définitions
\definecolor{ipssiVert}{HTML}{22C55E}       % À retenir
\definecolor{ipssiOrange}{HTML}{F97316}     % Attention / Piège
\definecolor{ipssiViolet}{HTML}{A855F7}     % Intuition
\definecolor{ipssiCyan}{HTML}{06B6D4}       % Exemples
\definecolor{ipssiGrisClair}{HTML}{F3F4F6}  % Fond code
\definecolor{ipssiGrisFonce}{HTML}{1F2937}  % Texte code
\definecolor{ipssiFondPage}{HTML}{FAFAFA}   % Fond léger page de garde
\definecolor{ipssiRouge}{HTML}{E74C3C}      % Classe 0 (salaire ≤ 45k€)
\definecolor{ipssiVertData}{HTML}{2ECC71}   % Classe 1 (salaire > 45k€)
\definecolor{ipssiBleuDroite}{HTML}{3498DB} % Droite de régression

% ============================================================================
% ENCADRÉS PÉDAGOGIQUES (tcolorbox)
% ============================================================================
\usepackage[most]{tcolorbox}

% Style de base partagé
\tcbset{
    ipssibase/.style={
        enhanced,
        breakable,
        left=4mm,
        right=4mm,
        top=3mm,
        bottom=3mm,
        boxrule=0pt,
        frame hidden,
        borderline west={3pt}{0pt}{#1},
        colback=#1!5,
        colframe=#1,
        fonttitle=\bfseries\sffamily,
        coltitle=#1!80!black,
        attach boxed title to top left={yshift=-2mm, xshift=4mm},
        boxed title style={
            colback=#1!15,
            colframe=#1,
            boxrule=0.5pt,
            arc=2pt,
        },
        arc=2pt,
        outer arc=2pt,
    }
}

% Icônes TikZ pour les encadrés (compatible pdflatex, pas d'emoji Unicode)
\newcommand{\icoDefinition}{%
    \tikz[baseline=-0.3em]{\draw[ipssiBleu, thick] (0,0) -- (0.3,0) -- (0.3,0.35) -- (0,0.35) -- cycle;
    \draw[ipssiBleu, thick] (0.05,0.1) -- (0.25,0.1); \draw[ipssiBleu, thick] (0.05,0.2) -- (0.2,0.2);}%
}
\newcommand{\icoRetenir}{%
    \tikz[baseline=-0.2em]{\draw[ipssiVert, thick, fill=ipssiVert!20] (0.15,0) -- (0.3,0.3) -- (0,0.3) -- cycle;}%
}
\newcommand{\icoAttention}{%
    \tikz[baseline=-0.2em]{\draw[ipssiOrange, thick] (0.15,0.35) -- (0,0) -- (0.3,0) -- cycle;
    \node[ipssiOrange, font=\tiny\bfseries] at (0.15,0.12) {!};}%
}
\newcommand{\icoIntuition}{%
    \tikz[baseline=-0.2em]{\draw[ipssiViolet, thick, fill=ipssiViolet!15]
    (0.15,0.35) -- (0.08,0.2) -- (0.08,0.15) -- (0.22,0.15) -- (0.22,0.2) -- cycle;
    \draw[ipssiViolet, thick] (0.11,0.12) -- (0.11,0.05); \draw[ipssiViolet, thick] (0.19,0.12) -- (0.19,0.05);}%
}
\newcommand{\icoExemple}{%
    \tikz[baseline=-0.2em]{\draw[ipssiCyan, thick] (0,0) rectangle (0.3,0.3);
    \draw[ipssiCyan, thick] (0,0.15) -- (0.3,0.15); \draw[ipssiCyan, thick] (0.15,0) -- (0.15,0.3);}%
}
\newcommand{\icoPython}{%
    \tikz[baseline=-0.2em]{\node[font=\small\bfseries\ttfamily, ipssiVert!70!black] at (0,0) {\textgreater\textgreater};}%
}

% Définition (bleu)
\newtcolorbox{definitionbox}[1][D\'{e}finition]{
    ipssibase=ipssiBleu,
    title={\icoDefinition\; #1},
}

% À retenir (vert)
\newtcolorbox{retenirbox}[1][\`{A} retenir]{
    ipssibase=ipssiVert,
    title={\icoRetenir\; #1},
}

% Attention / Piège (orange)
\newtcolorbox{attentionbox}[1][Attention]{
    ipssibase=ipssiOrange,
    title={\icoAttention\; #1},
}

% Intuition (violet)
\newtcolorbox{intuitionbox}[1][Intuition]{
    ipssibase=ipssiViolet,
    title={\icoIntuition\; #1},
}

% Exemple (cyan)
\newtcolorbox{exemplebox}[1][Exemple]{
    ipssibase=ipssiCyan,
    title={\icoExemple\; #1},
}

% Code Python (style terminal)
\newtcolorbox{pythonbox}[1][Code Python]{
    enhanced,
    breakable,
    left=4mm, right=4mm, top=3mm, bottom=3mm,
    boxrule=0.5pt,
    colback=ipssiGrisClair,
    colframe=ipssiGrisFonce!40,
    fonttitle=\bfseries\ttfamily\small,
    coltitle=ipssiGrisFonce,
    title={\icoPython\; #1},
    attach boxed title to top left={yshift=-2mm, xshift=4mm},
    boxed title style={colback=ipssiGrisClair, colframe=ipssiGrisFonce!40, boxrule=0.5pt, arc=2pt},
    arc=2pt, outer arc=2pt,
}

% ============================================================================
% CODE PYTHON (listings)
% ============================================================================
\usepackage{listings}

\lstdefinestyle{python}{
    language=Python,
    basicstyle=\ttfamily\small,
    keywordstyle=\color{ipssiBleu}\bfseries,
    stringstyle=\color{ipssiVert!70!black},
    commentstyle=\color{gray}\itshape,
    numberstyle=\tiny\color{gray},
    numbers=left,
    numbersep=8pt,
    backgroundcolor=\color{ipssiGrisClair},
    frame=none,
    rulecolor=\color{ipssiGrisFonce!20},
    breaklines=true,
    breakatwhitespace=true,
    tabsize=4,
    showstringspaces=false,
    captionpos=b,
    aboveskip=8pt,
    belowskip=8pt,
    xleftmargin=10pt,
    framexleftmargin=10pt,
    morekeywords={as, with, True, False, None, self, np, plt, import, from},
    literate=
        {é}{{\'e}}1 {è}{{\`e}}1 {ê}{{\^e}}1 {ë}{{\"e}}1
        {à}{{\`a}}1 {â}{{\^a}}1
        {ù}{{\`u}}1 {û}{{\^u}}1
        {ô}{{\^o}}1
        {î}{{\^i}}1 {ï}{{\"i}}1
        {ç}{{\c{c}}}1
        {É}{{\'E}}1 {È}{{\`E}}1
        {→}{{$\rightarrow$}}1
        {≈}{{$\approx$}}1
        {≤}{{$\leq$}}1
        {≥}{{$\geq$}}1
        {★}{{$\bigstar$}}1
        {ŷ}{{$\hat{y}$}}1
        {·}{{$\cdot$}}1,
}

\lstset{style=python}

% Style pour le code inline
\newcommand{\pyinline}[1]{\lstinline[style=python,basicstyle=\ttfamily\small]{#1}}

% ============================================================================
% GRAPHIQUES
% ============================================================================
\usepackage{tikz}
\usepackage{pgfplots}
\pgfplotsset{compat=1.18}
\usetikzlibrary{arrows.meta, calc, positioning, shapes.geometric, decorations.pathreplacing}

% ============================================================================
% TABLEAUX
% ============================================================================
\usepackage{booktabs}
\usepackage{tabularx}
\usepackage{multirow}

% ============================================================================
% LISTES
% ============================================================================
\usepackage{enumitem}
\setlist[itemize]{itemsep=2pt, parsep=0pt, topsep=4pt}
\setlist[enumerate]{itemsep=2pt, parsep=0pt, topsep=4pt}

% ============================================================================
% HYPERLIENS
% ============================================================================
\usepackage[
    colorlinks=true,
    linkcolor=ipssiBleu!80!black,
    urlcolor=ipssiBleu!80!black,
    citecolor=ipssiVert!80!black,
    bookmarks=true,
    bookmarksopen=true,
    pdfauthor={Mohamed EL AFRIT},
    pdftitle={Mathématiques appliquées à l'IA — IPSSI},
    pdfsubject={Formation Bachelor 2 — IPSSI 2025-2026},
]{hyperref}

% ============================================================================
% IMAGES
% ============================================================================
\usepackage{graphicx}

% ============================================================================
% VARIABLES PARAMÉTRABLES (jour et thème)
% ============================================================================
\newcommand{\leJour}{X}
\newcommand{\leTheme}{Thème}
\newcommand{\jourNumero}[1]{\renewcommand{\leJour}{#1}}
\newcommand{\jourTheme}[1]{\renewcommand{\leTheme}{#1}}

% ============================================================================
% EN-TÊTE ET PIED DE PAGE (fancyhdr)
% ============================================================================
\usepackage{fancyhdr}

\pagestyle{fancy}
\fancyhf{}

% En-tête
\fancyhead[L]{\small\sffamily\textcolor{gray}{IPSSI — Maths appliquées à l'IA}}
\fancyhead[R]{\small\sffamily\textcolor{gray}{Jour \leJour\ — \leTheme}}
\renewcommand{\headrulewidth}{0.4pt}

% Pied de page
\fancyfoot[L]{\footnotesize\textcolor{gray}{Mohamed EL AFRIT\enspace\textbullet\enspace\url{www.mohamedelafrit.com}\enspace\textbullet\enspace CC BY-NC-SA 4.0}}
\fancyfoot[C]{\footnotesize\textcolor{gray}{\thepage}}
\fancyfoot[R]{\footnotesize\textcolor{gray}{2025–2026}}
\renewcommand{\footrulewidth}{0.2pt}

% Style pour la page de garde (pas d'en-tête)
\fancypagestyle{pagegarde}{
    \fancyhf{}
    \renewcommand{\headrulewidth}{0pt}
    \fancyfoot[C]{\footnotesize\textcolor{gray}{\thepage}}
    \renewcommand{\footrulewidth}{0pt}
}

% ============================================================================
% PAGE DE GARDE — \pagegarde{titre}{jour}{theme}
% ============================================================================
\newcommand{\pagegarde}[3]{%
    \thispagestyle{pagegarde}
    \begin{center}
        \vspace*{1.5cm}

        % Logo / Bandeau institution
        {\Large\sffamily\bfseries\textcolor{ipssiBleu}{IPSSI}}\par
        \vspace{2pt}
        {\normalsize\sffamily 2\textsuperscript{e} année Bachelor Informatique}\par

        \vspace{0.8cm}
        {\color{ipssiBleu}\rule{0.6\textwidth}{1.5pt}}\par
        \vspace{0.8cm}

        % Titre du module
        {\LARGE\sffamily\bfseries Mathématiques appliquées\\[4pt] à l'Intelligence Artificielle}\par

        \vspace{1cm}

        % Titre du document
        \begin{tcolorbox}[
            enhanced,
            colback=ipssiBleu!5,
            colframe=ipssiBleu!50,
            boxrule=1pt,
            arc=4pt,
            width=0.85\textwidth,
            halign=center,
            top=10pt, bottom=10pt,
        ]
            {\huge\sffamily\bfseries\textcolor{ipssiBleu!80!black}{#1}}\par
            \vspace{6pt}
            {\Large\sffamily Jour #2 — #3}
        \end{tcolorbox}

        \vspace{1.5cm}

        % Informations auteur
        {\large\sffamily
            \textbf{Auteur :} Mohamed EL AFRIT\\[4pt]
            \url{www.mohamedelafrit.com}\\[8pt]
            \textbf{Année académique :} 2025–2026
        }\par

        \vspace{1.5cm}

        % Licence Creative Commons
        {\color{ipssiBleu}\rule{0.4\textwidth}{0.5pt}}\par
        \vspace{8pt}
        {\small\sffamily
            \textcopyright\ 2025 Mohamed EL AFRIT — IPSSI\\[2pt]
            Ce contenu est distribué sous licence\\[2pt]
            \textbf{Creative Commons BY-NC-SA 4.0}\\[2pt]
            \url{https://creativecommons.org/licenses/by-nc-sa/4.0/deed.fr}
        }\par
    \end{center}
    \newpage
}

% ============================================================================
% NIVEAUX D'EXERCICES
% ============================================================================

% Étoile compatible pdflatex
\newcommand{\starfull}[1]{\textcolor{#1}{$\bigstar$}}
\newcommand{\starempty}{\textcolor{gray!40}{$\bigstar$}}

% Fondamental (1 étoile)
\newcommand{\exofondamental}{%
    \starfull{ipssiVert}\starempty\starempty\enspace%
    {\small\sffamily\textcolor{ipssiVert}{\textbf{Fondamental}}}%
}

% Intermédiaire (2 étoiles)
\newcommand{\exointermediaire}{%
    \starfull{ipssiOrange}\starfull{ipssiOrange}\starempty\enspace%
    {\small\sffamily\textcolor{ipssiOrange}{\textbf{Interm\'{e}diaire}}}%
}

% Avancé (3 étoiles)
\newcommand{\exoavance}{%
    \starfull{ipssiRouge}\starfull{ipssiRouge}\starfull{ipssiRouge}\enspace%
    {\small\sffamily\textcolor{ipssiRouge}{\textbf{Avanc\'{e}}}}%
}

% Commande d'exercice numéroté
\newcounter{exocount}[section]
\newcommand{\exercice}[1]{%
    \stepcounter{exocount}%
    \vspace{8pt}%
    \noindent\textbf{\sffamily Exercice \theexocount}\enspace — \enspace #1\par
    \vspace{4pt}%
}

% ============================================================================
% COMMANDES UTILITAIRES
% ============================================================================

% Séparateur visuel entre sections
\newcommand{\separateur}{%
    \vspace{6pt}
    \begin{center}
        {\color{ipssiBleu!30}\rule{0.3\textwidth}{0.5pt}}
    \end{center}
    \vspace{6pt}
}

% Mot-clé en gras coloré
\newcommand{\motcle}[1]{\textbf{\textcolor{ipssiBleu!80!black}{#1}}}

% Formule encadrée importante
\newcommand{\formule}[1]{%
    \begin{center}
        \fcolorbox{ipssiBleu!50}{ipssiBleu!5}{%
            \quad$\displaystyle #1$\quad%
        }
    \end{center}
}

% ============================================================================
% FIN DU PRÉAMBULE
% ============================================================================
          % Importer ce préambule
%   \jourNumero{1}                  % Définir le numéro du jour (1, 2, 3 ou 4)
%   \jourTheme{Données et représentation mathématique}
%   \begin{document}
%   \pagegarde{Fiche de synthèse}{1}{Données et représentation mathématique}
%   ...
%   \end{document}
% ============================================================================

% --- Classe de document ---
\documentclass[11pt,a4paper]{article}

% ============================================================================
% ENCODAGE ET LANGUE
% ============================================================================
\usepackage[utf8]{inputenc}
\usepackage[T1]{fontenc}
\usepackage[french]{babel}

% ============================================================================
% MATHÉMATIQUES
% ============================================================================
\usepackage{amsmath,amssymb,amsfonts,mathtools}

% Commandes mathématiques personnalisées (conventions du cours)
\newcommand{\vx}{\mathbf{x}}          % vecteur x
\newcommand{\vw}{\mathbf{w}}          % vecteur w (poids)
\newcommand{\vy}{\mathbf{y}}          % vecteur y
\newcommand{\mX}{\mathbf{X}}          % matrice X
\newcommand{\mW}{\mathbf{W}}          % matrice W
\newcommand{\yhat}{\hat{y}}           % prédiction
\newcommand{\pscal}[2]{\langle #1, #2 \rangle} % produit scalaire
\newcommand{\gradJ}{\nabla J(\boldsymbol{\theta})} % gradient de J

% ============================================================================
% MISE EN PAGE
% ============================================================================
\usepackage[margin=2.5cm]{geometry}
\usepackage{setspace}
\setstretch{1.15}                      % interligne légèrement aéré

% ============================================================================
% COULEURS
% ============================================================================
\usepackage[dvipsnames,svgnames,x11names]{xcolor}

% Palette principale du projet
\definecolor{ipssiBleu}{HTML}{3B82F6}       % Définitions
\definecolor{ipssiVert}{HTML}{22C55E}       % À retenir
\definecolor{ipssiOrange}{HTML}{F97316}     % Attention / Piège
\definecolor{ipssiViolet}{HTML}{A855F7}     % Intuition
\definecolor{ipssiCyan}{HTML}{06B6D4}       % Exemples
\definecolor{ipssiGrisClair}{HTML}{F3F4F6}  % Fond code
\definecolor{ipssiGrisFonce}{HTML}{1F2937}  % Texte code
\definecolor{ipssiFondPage}{HTML}{FAFAFA}   % Fond léger page de garde
\definecolor{ipssiRouge}{HTML}{E74C3C}      % Classe 0 (salaire ≤ 45k€)
\definecolor{ipssiVertData}{HTML}{2ECC71}   % Classe 1 (salaire > 45k€)
\definecolor{ipssiBleuDroite}{HTML}{3498DB} % Droite de régression

% ============================================================================
% ENCADRÉS PÉDAGOGIQUES (tcolorbox)
% ============================================================================
\usepackage[most]{tcolorbox}

% Style de base partagé
\tcbset{
    ipssibase/.style={
        enhanced,
        breakable,
        left=4mm,
        right=4mm,
        top=3mm,
        bottom=3mm,
        boxrule=0pt,
        frame hidden,
        borderline west={3pt}{0pt}{#1},
        colback=#1!5,
        colframe=#1,
        fonttitle=\bfseries\sffamily,
        coltitle=#1!80!black,
        attach boxed title to top left={yshift=-2mm, xshift=4mm},
        boxed title style={
            colback=#1!15,
            colframe=#1,
            boxrule=0.5pt,
            arc=2pt,
        },
        arc=2pt,
        outer arc=2pt,
    }
}

% Icônes TikZ pour les encadrés (compatible pdflatex, pas d'emoji Unicode)
\newcommand{\icoDefinition}{%
    \tikz[baseline=-0.3em]{\draw[ipssiBleu, thick] (0,0) -- (0.3,0) -- (0.3,0.35) -- (0,0.35) -- cycle;
    \draw[ipssiBleu, thick] (0.05,0.1) -- (0.25,0.1); \draw[ipssiBleu, thick] (0.05,0.2) -- (0.2,0.2);}%
}
\newcommand{\icoRetenir}{%
    \tikz[baseline=-0.2em]{\draw[ipssiVert, thick, fill=ipssiVert!20] (0.15,0) -- (0.3,0.3) -- (0,0.3) -- cycle;}%
}
\newcommand{\icoAttention}{%
    \tikz[baseline=-0.2em]{\draw[ipssiOrange, thick] (0.15,0.35) -- (0,0) -- (0.3,0) -- cycle;
    \node[ipssiOrange, font=\tiny\bfseries] at (0.15,0.12) {!};}%
}
\newcommand{\icoIntuition}{%
    \tikz[baseline=-0.2em]{\draw[ipssiViolet, thick, fill=ipssiViolet!15]
    (0.15,0.35) -- (0.08,0.2) -- (0.08,0.15) -- (0.22,0.15) -- (0.22,0.2) -- cycle;
    \draw[ipssiViolet, thick] (0.11,0.12) -- (0.11,0.05); \draw[ipssiViolet, thick] (0.19,0.12) -- (0.19,0.05);}%
}
\newcommand{\icoExemple}{%
    \tikz[baseline=-0.2em]{\draw[ipssiCyan, thick] (0,0) rectangle (0.3,0.3);
    \draw[ipssiCyan, thick] (0,0.15) -- (0.3,0.15); \draw[ipssiCyan, thick] (0.15,0) -- (0.15,0.3);}%
}
\newcommand{\icoPython}{%
    \tikz[baseline=-0.2em]{\node[font=\small\bfseries\ttfamily, ipssiVert!70!black] at (0,0) {\textgreater\textgreater};}%
}

% Définition (bleu)
\newtcolorbox{definitionbox}[1][D\'{e}finition]{
    ipssibase=ipssiBleu,
    title={\icoDefinition\; #1},
}

% À retenir (vert)
\newtcolorbox{retenirbox}[1][\`{A} retenir]{
    ipssibase=ipssiVert,
    title={\icoRetenir\; #1},
}

% Attention / Piège (orange)
\newtcolorbox{attentionbox}[1][Attention]{
    ipssibase=ipssiOrange,
    title={\icoAttention\; #1},
}

% Intuition (violet)
\newtcolorbox{intuitionbox}[1][Intuition]{
    ipssibase=ipssiViolet,
    title={\icoIntuition\; #1},
}

% Exemple (cyan)
\newtcolorbox{exemplebox}[1][Exemple]{
    ipssibase=ipssiCyan,
    title={\icoExemple\; #1},
}

% Code Python (style terminal)
\newtcolorbox{pythonbox}[1][Code Python]{
    enhanced,
    breakable,
    left=4mm, right=4mm, top=3mm, bottom=3mm,
    boxrule=0.5pt,
    colback=ipssiGrisClair,
    colframe=ipssiGrisFonce!40,
    fonttitle=\bfseries\ttfamily\small,
    coltitle=ipssiGrisFonce,
    title={\icoPython\; #1},
    attach boxed title to top left={yshift=-2mm, xshift=4mm},
    boxed title style={colback=ipssiGrisClair, colframe=ipssiGrisFonce!40, boxrule=0.5pt, arc=2pt},
    arc=2pt, outer arc=2pt,
}

% ============================================================================
% CODE PYTHON (listings)
% ============================================================================
\usepackage{listings}

\lstdefinestyle{python}{
    language=Python,
    basicstyle=\ttfamily\small,
    keywordstyle=\color{ipssiBleu}\bfseries,
    stringstyle=\color{ipssiVert!70!black},
    commentstyle=\color{gray}\itshape,
    numberstyle=\tiny\color{gray},
    numbers=left,
    numbersep=8pt,
    backgroundcolor=\color{ipssiGrisClair},
    frame=none,
    rulecolor=\color{ipssiGrisFonce!20},
    breaklines=true,
    breakatwhitespace=true,
    tabsize=4,
    showstringspaces=false,
    captionpos=b,
    aboveskip=8pt,
    belowskip=8pt,
    xleftmargin=10pt,
    framexleftmargin=10pt,
    morekeywords={as, with, True, False, None, self, np, plt, import, from},
    literate=
        {é}{{\'e}}1 {è}{{\`e}}1 {ê}{{\^e}}1 {ë}{{\"e}}1
        {à}{{\`a}}1 {â}{{\^a}}1
        {ù}{{\`u}}1 {û}{{\^u}}1
        {ô}{{\^o}}1
        {î}{{\^i}}1 {ï}{{\"i}}1
        {ç}{{\c{c}}}1
        {É}{{\'E}}1 {È}{{\`E}}1
        {→}{{$\rightarrow$}}1
        {≈}{{$\approx$}}1
        {≤}{{$\leq$}}1
        {≥}{{$\geq$}}1
        {★}{{$\bigstar$}}1
        {ŷ}{{$\hat{y}$}}1
        {·}{{$\cdot$}}1,
}

\lstset{style=python}

% Style pour le code inline
\newcommand{\pyinline}[1]{\lstinline[style=python,basicstyle=\ttfamily\small]{#1}}

% ============================================================================
% GRAPHIQUES
% ============================================================================
\usepackage{tikz}
\usepackage{pgfplots}
\pgfplotsset{compat=1.18}
\usetikzlibrary{arrows.meta, calc, positioning, shapes.geometric, decorations.pathreplacing}

% ============================================================================
% TABLEAUX
% ============================================================================
\usepackage{booktabs}
\usepackage{tabularx}
\usepackage{multirow}

% ============================================================================
% LISTES
% ============================================================================
\usepackage{enumitem}
\setlist[itemize]{itemsep=2pt, parsep=0pt, topsep=4pt}
\setlist[enumerate]{itemsep=2pt, parsep=0pt, topsep=4pt}

% ============================================================================
% HYPERLIENS
% ============================================================================
\usepackage[
    colorlinks=true,
    linkcolor=ipssiBleu!80!black,
    urlcolor=ipssiBleu!80!black,
    citecolor=ipssiVert!80!black,
    bookmarks=true,
    bookmarksopen=true,
    pdfauthor={Mohamed EL AFRIT},
    pdftitle={Mathématiques appliquées à l'IA — IPSSI},
    pdfsubject={Formation Bachelor 2 — IPSSI 2025-2026},
]{hyperref}

% ============================================================================
% IMAGES
% ============================================================================
\usepackage{graphicx}

% ============================================================================
% VARIABLES PARAMÉTRABLES (jour et thème)
% ============================================================================
\newcommand{\leJour}{X}
\newcommand{\leTheme}{Thème}
\newcommand{\jourNumero}[1]{\renewcommand{\leJour}{#1}}
\newcommand{\jourTheme}[1]{\renewcommand{\leTheme}{#1}}

% ============================================================================
% EN-TÊTE ET PIED DE PAGE (fancyhdr)
% ============================================================================
\usepackage{fancyhdr}

\pagestyle{fancy}
\fancyhf{}

% En-tête
\fancyhead[L]{\small\sffamily\textcolor{gray}{IPSSI — Maths appliquées à l'IA}}
\fancyhead[R]{\small\sffamily\textcolor{gray}{Jour \leJour\ — \leTheme}}
\renewcommand{\headrulewidth}{0.4pt}

% Pied de page
\fancyfoot[L]{\footnotesize\textcolor{gray}{Mohamed EL AFRIT\enspace\textbullet\enspace\url{www.mohamedelafrit.com}\enspace\textbullet\enspace CC BY-NC-SA 4.0}}
\fancyfoot[C]{\footnotesize\textcolor{gray}{\thepage}}
\fancyfoot[R]{\footnotesize\textcolor{gray}{2025–2026}}
\renewcommand{\footrulewidth}{0.2pt}

% Style pour la page de garde (pas d'en-tête)
\fancypagestyle{pagegarde}{
    \fancyhf{}
    \renewcommand{\headrulewidth}{0pt}
    \fancyfoot[C]{\footnotesize\textcolor{gray}{\thepage}}
    \renewcommand{\footrulewidth}{0pt}
}

% ============================================================================
% PAGE DE GARDE — \pagegarde{titre}{jour}{theme}
% ============================================================================
\newcommand{\pagegarde}[3]{%
    \thispagestyle{pagegarde}
    \begin{center}
        \vspace*{1.5cm}

        % Logo / Bandeau institution
        {\Large\sffamily\bfseries\textcolor{ipssiBleu}{IPSSI}}\par
        \vspace{2pt}
        {\normalsize\sffamily 2\textsuperscript{e} année Bachelor Informatique}\par

        \vspace{0.8cm}
        {\color{ipssiBleu}\rule{0.6\textwidth}{1.5pt}}\par
        \vspace{0.8cm}

        % Titre du module
        {\LARGE\sffamily\bfseries Mathématiques appliquées\\[4pt] à l'Intelligence Artificielle}\par

        \vspace{1cm}

        % Titre du document
        \begin{tcolorbox}[
            enhanced,
            colback=ipssiBleu!5,
            colframe=ipssiBleu!50,
            boxrule=1pt,
            arc=4pt,
            width=0.85\textwidth,
            halign=center,
            top=10pt, bottom=10pt,
        ]
            {\huge\sffamily\bfseries\textcolor{ipssiBleu!80!black}{#1}}\par
            \vspace{6pt}
            {\Large\sffamily Jour #2 — #3}
        \end{tcolorbox}

        \vspace{1.5cm}

        % Informations auteur
        {\large\sffamily
            \textbf{Auteur :} Mohamed EL AFRIT\\[4pt]
            \url{www.mohamedelafrit.com}\\[8pt]
            \textbf{Année académique :} 2025–2026
        }\par

        \vspace{1.5cm}

        % Licence Creative Commons
        {\color{ipssiBleu}\rule{0.4\textwidth}{0.5pt}}\par
        \vspace{8pt}
        {\small\sffamily
            \textcopyright\ 2025 Mohamed EL AFRIT — IPSSI\\[2pt]
            Ce contenu est distribué sous licence\\[2pt]
            \textbf{Creative Commons BY-NC-SA 4.0}\\[2pt]
            \url{https://creativecommons.org/licenses/by-nc-sa/4.0/deed.fr}
        }\par
    \end{center}
    \newpage
}

% ============================================================================
% NIVEAUX D'EXERCICES
% ============================================================================

% Étoile compatible pdflatex
\newcommand{\starfull}[1]{\textcolor{#1}{$\bigstar$}}
\newcommand{\starempty}{\textcolor{gray!40}{$\bigstar$}}

% Fondamental (1 étoile)
\newcommand{\exofondamental}{%
    \starfull{ipssiVert}\starempty\starempty\enspace%
    {\small\sffamily\textcolor{ipssiVert}{\textbf{Fondamental}}}%
}

% Intermédiaire (2 étoiles)
\newcommand{\exointermediaire}{%
    \starfull{ipssiOrange}\starfull{ipssiOrange}\starempty\enspace%
    {\small\sffamily\textcolor{ipssiOrange}{\textbf{Interm\'{e}diaire}}}%
}

% Avancé (3 étoiles)
\newcommand{\exoavance}{%
    \starfull{ipssiRouge}\starfull{ipssiRouge}\starfull{ipssiRouge}\enspace%
    {\small\sffamily\textcolor{ipssiRouge}{\textbf{Avanc\'{e}}}}%
}

% Commande d'exercice numéroté
\newcounter{exocount}[section]
\newcommand{\exercice}[1]{%
    \stepcounter{exocount}%
    \vspace{8pt}%
    \noindent\textbf{\sffamily Exercice \theexocount}\enspace — \enspace #1\par
    \vspace{4pt}%
}

% ============================================================================
% COMMANDES UTILITAIRES
% ============================================================================

% Séparateur visuel entre sections
\newcommand{\separateur}{%
    \vspace{6pt}
    \begin{center}
        {\color{ipssiBleu!30}\rule{0.3\textwidth}{0.5pt}}
    \end{center}
    \vspace{6pt}
}

% Mot-clé en gras coloré
\newcommand{\motcle}[1]{\textbf{\textcolor{ipssiBleu!80!black}{#1}}}

% Formule encadrée importante
\newcommand{\formule}[1]{%
    \begin{center}
        \fcolorbox{ipssiBleu!50}{ipssiBleu!5}{%
            \quad$\displaystyle #1$\quad%
        }
    \end{center}
}

% ============================================================================
% FIN DU PRÉAMBULE
% ============================================================================
          % Importer ce préambule
%   \jourNumero{1}                  % Définir le numéro du jour (1, 2, 3 ou 4)
%   \jourTheme{Données et représentation mathématique}
%   \begin{document}
%   \pagegarde{Fiche de synthèse}{1}{Données et représentation mathématique}
%   ...
%   \end{document}
% ============================================================================

% --- Classe de document ---
\documentclass[11pt,a4paper]{article}

% ============================================================================
% ENCODAGE ET LANGUE
% ============================================================================
\usepackage[utf8]{inputenc}
\usepackage[T1]{fontenc}
\usepackage[french]{babel}

% ============================================================================
% MATHÉMATIQUES
% ============================================================================
\usepackage{amsmath,amssymb,amsfonts,mathtools}

% Commandes mathématiques personnalisées (conventions du cours)
\newcommand{\vx}{\mathbf{x}}          % vecteur x
\newcommand{\vw}{\mathbf{w}}          % vecteur w (poids)
\newcommand{\vy}{\mathbf{y}}          % vecteur y
\newcommand{\mX}{\mathbf{X}}          % matrice X
\newcommand{\mW}{\mathbf{W}}          % matrice W
\newcommand{\yhat}{\hat{y}}           % prédiction
\newcommand{\pscal}[2]{\langle #1, #2 \rangle} % produit scalaire
\newcommand{\gradJ}{\nabla J(\boldsymbol{\theta})} % gradient de J

% ============================================================================
% MISE EN PAGE
% ============================================================================
\usepackage[margin=2.5cm]{geometry}
\usepackage{setspace}
\setstretch{1.15}                      % interligne légèrement aéré

% ============================================================================
% COULEURS
% ============================================================================
\usepackage[dvipsnames,svgnames,x11names]{xcolor}

% Palette principale du projet
\definecolor{ipssiBleu}{HTML}{3B82F6}       % Définitions
\definecolor{ipssiVert}{HTML}{22C55E}       % À retenir
\definecolor{ipssiOrange}{HTML}{F97316}     % Attention / Piège
\definecolor{ipssiViolet}{HTML}{A855F7}     % Intuition
\definecolor{ipssiCyan}{HTML}{06B6D4}       % Exemples
\definecolor{ipssiGrisClair}{HTML}{F3F4F6}  % Fond code
\definecolor{ipssiGrisFonce}{HTML}{1F2937}  % Texte code
\definecolor{ipssiFondPage}{HTML}{FAFAFA}   % Fond léger page de garde
\definecolor{ipssiRouge}{HTML}{E74C3C}      % Classe 0 (salaire ≤ 45k€)
\definecolor{ipssiVertData}{HTML}{2ECC71}   % Classe 1 (salaire > 45k€)
\definecolor{ipssiBleuDroite}{HTML}{3498DB} % Droite de régression

% ============================================================================
% ENCADRÉS PÉDAGOGIQUES (tcolorbox)
% ============================================================================
\usepackage[most]{tcolorbox}

% Style de base partagé
\tcbset{
    ipssibase/.style={
        enhanced,
        breakable,
        left=4mm,
        right=4mm,
        top=3mm,
        bottom=3mm,
        boxrule=0pt,
        frame hidden,
        borderline west={3pt}{0pt}{#1},
        colback=#1!5,
        colframe=#1,
        fonttitle=\bfseries\sffamily,
        coltitle=#1!80!black,
        attach boxed title to top left={yshift=-2mm, xshift=4mm},
        boxed title style={
            colback=#1!15,
            colframe=#1,
            boxrule=0.5pt,
            arc=2pt,
        },
        arc=2pt,
        outer arc=2pt,
    }
}

% Icônes TikZ pour les encadrés (compatible pdflatex, pas d'emoji Unicode)
\newcommand{\icoDefinition}{%
    \tikz[baseline=-0.3em]{\draw[ipssiBleu, thick] (0,0) -- (0.3,0) -- (0.3,0.35) -- (0,0.35) -- cycle;
    \draw[ipssiBleu, thick] (0.05,0.1) -- (0.25,0.1); \draw[ipssiBleu, thick] (0.05,0.2) -- (0.2,0.2);}%
}
\newcommand{\icoRetenir}{%
    \tikz[baseline=-0.2em]{\draw[ipssiVert, thick, fill=ipssiVert!20] (0.15,0) -- (0.3,0.3) -- (0,0.3) -- cycle;}%
}
\newcommand{\icoAttention}{%
    \tikz[baseline=-0.2em]{\draw[ipssiOrange, thick] (0.15,0.35) -- (0,0) -- (0.3,0) -- cycle;
    \node[ipssiOrange, font=\tiny\bfseries] at (0.15,0.12) {!};}%
}
\newcommand{\icoIntuition}{%
    \tikz[baseline=-0.2em]{\draw[ipssiViolet, thick, fill=ipssiViolet!15]
    (0.15,0.35) -- (0.08,0.2) -- (0.08,0.15) -- (0.22,0.15) -- (0.22,0.2) -- cycle;
    \draw[ipssiViolet, thick] (0.11,0.12) -- (0.11,0.05); \draw[ipssiViolet, thick] (0.19,0.12) -- (0.19,0.05);}%
}
\newcommand{\icoExemple}{%
    \tikz[baseline=-0.2em]{\draw[ipssiCyan, thick] (0,0) rectangle (0.3,0.3);
    \draw[ipssiCyan, thick] (0,0.15) -- (0.3,0.15); \draw[ipssiCyan, thick] (0.15,0) -- (0.15,0.3);}%
}
\newcommand{\icoPython}{%
    \tikz[baseline=-0.2em]{\node[font=\small\bfseries\ttfamily, ipssiVert!70!black] at (0,0) {\textgreater\textgreater};}%
}

% Définition (bleu)
\newtcolorbox{definitionbox}[1][D\'{e}finition]{
    ipssibase=ipssiBleu,
    title={\icoDefinition\; #1},
}

% À retenir (vert)
\newtcolorbox{retenirbox}[1][\`{A} retenir]{
    ipssibase=ipssiVert,
    title={\icoRetenir\; #1},
}

% Attention / Piège (orange)
\newtcolorbox{attentionbox}[1][Attention]{
    ipssibase=ipssiOrange,
    title={\icoAttention\; #1},
}

% Intuition (violet)
\newtcolorbox{intuitionbox}[1][Intuition]{
    ipssibase=ipssiViolet,
    title={\icoIntuition\; #1},
}

% Exemple (cyan)
\newtcolorbox{exemplebox}[1][Exemple]{
    ipssibase=ipssiCyan,
    title={\icoExemple\; #1},
}

% Code Python (style terminal)
\newtcolorbox{pythonbox}[1][Code Python]{
    enhanced,
    breakable,
    left=4mm, right=4mm, top=3mm, bottom=3mm,
    boxrule=0.5pt,
    colback=ipssiGrisClair,
    colframe=ipssiGrisFonce!40,
    fonttitle=\bfseries\ttfamily\small,
    coltitle=ipssiGrisFonce,
    title={\icoPython\; #1},
    attach boxed title to top left={yshift=-2mm, xshift=4mm},
    boxed title style={colback=ipssiGrisClair, colframe=ipssiGrisFonce!40, boxrule=0.5pt, arc=2pt},
    arc=2pt, outer arc=2pt,
}

% ============================================================================
% CODE PYTHON (listings)
% ============================================================================
\usepackage{listings}

\lstdefinestyle{python}{
    language=Python,
    basicstyle=\ttfamily\small,
    keywordstyle=\color{ipssiBleu}\bfseries,
    stringstyle=\color{ipssiVert!70!black},
    commentstyle=\color{gray}\itshape,
    numberstyle=\tiny\color{gray},
    numbers=left,
    numbersep=8pt,
    backgroundcolor=\color{ipssiGrisClair},
    frame=none,
    rulecolor=\color{ipssiGrisFonce!20},
    breaklines=true,
    breakatwhitespace=true,
    tabsize=4,
    showstringspaces=false,
    captionpos=b,
    aboveskip=8pt,
    belowskip=8pt,
    xleftmargin=10pt,
    framexleftmargin=10pt,
    morekeywords={as, with, True, False, None, self, np, plt, import, from},
    literate=
        {é}{{\'e}}1 {è}{{\`e}}1 {ê}{{\^e}}1 {ë}{{\"e}}1
        {à}{{\`a}}1 {â}{{\^a}}1
        {ù}{{\`u}}1 {û}{{\^u}}1
        {ô}{{\^o}}1
        {î}{{\^i}}1 {ï}{{\"i}}1
        {ç}{{\c{c}}}1
        {É}{{\'E}}1 {È}{{\`E}}1
        {→}{{$\rightarrow$}}1
        {≈}{{$\approx$}}1
        {≤}{{$\leq$}}1
        {≥}{{$\geq$}}1
        {★}{{$\bigstar$}}1
        {ŷ}{{$\hat{y}$}}1
        {·}{{$\cdot$}}1,
}

\lstset{style=python}

% Style pour le code inline
\newcommand{\pyinline}[1]{\lstinline[style=python,basicstyle=\ttfamily\small]{#1}}

% ============================================================================
% GRAPHIQUES
% ============================================================================
\usepackage{tikz}
\usepackage{pgfplots}
\pgfplotsset{compat=1.18}
\usetikzlibrary{arrows.meta, calc, positioning, shapes.geometric, decorations.pathreplacing}

% ============================================================================
% TABLEAUX
% ============================================================================
\usepackage{booktabs}
\usepackage{tabularx}
\usepackage{multirow}

% ============================================================================
% LISTES
% ============================================================================
\usepackage{enumitem}
\setlist[itemize]{itemsep=2pt, parsep=0pt, topsep=4pt}
\setlist[enumerate]{itemsep=2pt, parsep=0pt, topsep=4pt}

% ============================================================================
% HYPERLIENS
% ============================================================================
\usepackage[
    colorlinks=true,
    linkcolor=ipssiBleu!80!black,
    urlcolor=ipssiBleu!80!black,
    citecolor=ipssiVert!80!black,
    bookmarks=true,
    bookmarksopen=true,
    pdfauthor={Mohamed EL AFRIT},
    pdftitle={Mathématiques appliquées à l'IA — IPSSI},
    pdfsubject={Formation Bachelor 2 — IPSSI 2025-2026},
]{hyperref}

% ============================================================================
% IMAGES
% ============================================================================
\usepackage{graphicx}

% ============================================================================
% VARIABLES PARAMÉTRABLES (jour et thème)
% ============================================================================
\newcommand{\leJour}{X}
\newcommand{\leTheme}{Thème}
\newcommand{\jourNumero}[1]{\renewcommand{\leJour}{#1}}
\newcommand{\jourTheme}[1]{\renewcommand{\leTheme}{#1}}

% ============================================================================
% EN-TÊTE ET PIED DE PAGE (fancyhdr)
% ============================================================================
\usepackage{fancyhdr}

\pagestyle{fancy}
\fancyhf{}

% En-tête
\fancyhead[L]{\small\sffamily\textcolor{gray}{IPSSI — Maths appliquées à l'IA}}
\fancyhead[R]{\small\sffamily\textcolor{gray}{Jour \leJour\ — \leTheme}}
\renewcommand{\headrulewidth}{0.4pt}

% Pied de page
\fancyfoot[L]{\footnotesize\textcolor{gray}{Mohamed EL AFRIT\enspace\textbullet\enspace\url{www.mohamedelafrit.com}\enspace\textbullet\enspace CC BY-NC-SA 4.0}}
\fancyfoot[C]{\footnotesize\textcolor{gray}{\thepage}}
\fancyfoot[R]{\footnotesize\textcolor{gray}{2025–2026}}
\renewcommand{\footrulewidth}{0.2pt}

% Style pour la page de garde (pas d'en-tête)
\fancypagestyle{pagegarde}{
    \fancyhf{}
    \renewcommand{\headrulewidth}{0pt}
    \fancyfoot[C]{\footnotesize\textcolor{gray}{\thepage}}
    \renewcommand{\footrulewidth}{0pt}
}

% ============================================================================
% PAGE DE GARDE — \pagegarde{titre}{jour}{theme}
% ============================================================================
\newcommand{\pagegarde}[3]{%
    \thispagestyle{pagegarde}
    \begin{center}
        \vspace*{1.5cm}

        % Logo / Bandeau institution
        {\Large\sffamily\bfseries\textcolor{ipssiBleu}{IPSSI}}\par
        \vspace{2pt}
        {\normalsize\sffamily 2\textsuperscript{e} année Bachelor Informatique}\par

        \vspace{0.8cm}
        {\color{ipssiBleu}\rule{0.6\textwidth}{1.5pt}}\par
        \vspace{0.8cm}

        % Titre du module
        {\LARGE\sffamily\bfseries Mathématiques appliquées\\[4pt] à l'Intelligence Artificielle}\par

        \vspace{1cm}

        % Titre du document
        \begin{tcolorbox}[
            enhanced,
            colback=ipssiBleu!5,
            colframe=ipssiBleu!50,
            boxrule=1pt,
            arc=4pt,
            width=0.85\textwidth,
            halign=center,
            top=10pt, bottom=10pt,
        ]
            {\huge\sffamily\bfseries\textcolor{ipssiBleu!80!black}{#1}}\par
            \vspace{6pt}
            {\Large\sffamily Jour #2 — #3}
        \end{tcolorbox}

        \vspace{1.5cm}

        % Informations auteur
        {\large\sffamily
            \textbf{Auteur :} Mohamed EL AFRIT\\[4pt]
            \url{www.mohamedelafrit.com}\\[8pt]
            \textbf{Année académique :} 2025–2026
        }\par

        \vspace{1.5cm}

        % Licence Creative Commons
        {\color{ipssiBleu}\rule{0.4\textwidth}{0.5pt}}\par
        \vspace{8pt}
        {\small\sffamily
            \textcopyright\ 2025 Mohamed EL AFRIT — IPSSI\\[2pt]
            Ce contenu est distribué sous licence\\[2pt]
            \textbf{Creative Commons BY-NC-SA 4.0}\\[2pt]
            \url{https://creativecommons.org/licenses/by-nc-sa/4.0/deed.fr}
        }\par
    \end{center}
    \newpage
}

% ============================================================================
% NIVEAUX D'EXERCICES
% ============================================================================

% Étoile compatible pdflatex
\newcommand{\starfull}[1]{\textcolor{#1}{$\bigstar$}}
\newcommand{\starempty}{\textcolor{gray!40}{$\bigstar$}}

% Fondamental (1 étoile)
\newcommand{\exofondamental}{%
    \starfull{ipssiVert}\starempty\starempty\enspace%
    {\small\sffamily\textcolor{ipssiVert}{\textbf{Fondamental}}}%
}

% Intermédiaire (2 étoiles)
\newcommand{\exointermediaire}{%
    \starfull{ipssiOrange}\starfull{ipssiOrange}\starempty\enspace%
    {\small\sffamily\textcolor{ipssiOrange}{\textbf{Interm\'{e}diaire}}}%
}

% Avancé (3 étoiles)
\newcommand{\exoavance}{%
    \starfull{ipssiRouge}\starfull{ipssiRouge}\starfull{ipssiRouge}\enspace%
    {\small\sffamily\textcolor{ipssiRouge}{\textbf{Avanc\'{e}}}}%
}

% Commande d'exercice numéroté
\newcounter{exocount}[section]
\newcommand{\exercice}[1]{%
    \stepcounter{exocount}%
    \vspace{8pt}%
    \noindent\textbf{\sffamily Exercice \theexocount}\enspace — \enspace #1\par
    \vspace{4pt}%
}

% ============================================================================
% COMMANDES UTILITAIRES
% ============================================================================

% Séparateur visuel entre sections
\newcommand{\separateur}{%
    \vspace{6pt}
    \begin{center}
        {\color{ipssiBleu!30}\rule{0.3\textwidth}{0.5pt}}
    \end{center}
    \vspace{6pt}
}

% Mot-clé en gras coloré
\newcommand{\motcle}[1]{\textbf{\textcolor{ipssiBleu!80!black}{#1}}}

% Formule encadrée importante
\newcommand{\formule}[1]{%
    \begin{center}
        \fcolorbox{ipssiBleu!50}{ipssiBleu!5}{%
            \quad$\displaystyle #1$\quad%
        }
    \end{center}
}

% ============================================================================
% FIN DU PRÉAMBULE
% ============================================================================


\jourNumero{1}
\jourTheme{Repr\'{e}sentation}

\begin{document}

\pagegarde{Corrections d\'{e}taill\'{e}es}{1}{Donn\'{e}es et repr\'{e}sentation math\'{e}matique}

\setstretch{1.06}
\setlength{\parskip}{2pt}
\setlength{\abovedisplayskip}{5pt}
\setlength{\belowdisplayskip}{5pt}

% ======================================================================
%  RAPPEL DATASET
% ======================================================================
\begin{tcolorbox}[
    enhanced, colback=ipssiCyan!5, colframe=ipssiCyan!50,
    boxrule=0.8pt, arc=4pt,
    title={\sffamily\bfseries Rappel --- Dataset fil rouge},
    coltitle=white, colbacktitle=ipssiCyan!70,
]
\small
\begin{center}
\renewcommand{\arraystretch}{1.1}
\begin{tabular}{c|ccccc|c}
    \toprule
    \textbf{\#} & \textbf{Exp} & \textbf{Lang} & \textbf{\'{E}tudes} & \textbf{Entrep.} & \textbf{Remote} & \textbf{Salaire} \\
    \midrule
    0 & 0 & 2 & 3 & 150  & 80  & 28.0 \\
    1 & 1 & 1 & 2 & 50   & 100 & 26.5 \\
    2 & 1 & 3 & 4 & 800  & 40  & 35.0 \\
    3 & 2 & 2 & 3 & 200  & 60  & 32.0 \\
    4 & 2 & 4 & 4 & 3000 & 20  & 38.0 \\
    5 & 0 & 1 & 1 & 30   & 50  & 25.0 \\
    \bottomrule
\end{tabular}
\end{center}
\end{tcolorbox}

% ======================================================================
% EXERCICE 1
% ======================================================================
\section*{Exercice 1 \hfill \exofondamental}
\addcontentsline{toc}{section}{Exercice 1 --- Vecteurs et dimensions}
\textit{Rappel : \'{e}crire des vecteurs de features et v\'{e}rifier les dimensions.}

\separateur

\textbf{a) Vecteur $\vx_0$ du d\'{e}veloppeur \#0.}

Le d\'{e}veloppeur \#0 a les features : exp=0, lang=2, \'{e}tudes=3, entrep.=150, remote=80.

\formule{
    \vx_0 = \begin{pmatrix} 0 \\ 2 \\ 3 \\ 150 \\ 80 \end{pmatrix} \in \mathbb{R}^5
    \qquad \Longrightarrow \qquad p = 5
}

\textbf{b) Vecteur $\vx_4$ et sa transpos\'{e}e.}

\[
    \vx_4 = \begin{pmatrix} 2 \\ 4 \\ 4 \\ 3000 \\ 20 \end{pmatrix}
    \qquad
    \vx_4^T = \begin{pmatrix} 2 & 4 & 4 & 3000 & 20 \end{pmatrix}
\]

\textbf{c) Nouveau d\'{e}veloppeur} (6 ans, 4 langages, Bac+5, 1200 employ\'{e}s, 50\% remote) :
\[
    \vx = \begin{pmatrix} 6 \\ 4 \\ 4 \\ 1200 \\ 50 \end{pmatrix} \in \mathbb{R}^5
\]

\textbf{d) Addition de vecteurs.}
$\vx_0 + \vx_4$ : \textbf{oui}, car les deux sont dans $\mathbb{R}^5$ (m\^{e}me dimension).
\[
    \vx_0 + \vx_4 = \begin{pmatrix} 0+2 \\ 2+4 \\ 3+4 \\ 150+3000 \\ 80+20 \end{pmatrix}
    = \begin{pmatrix} 2 \\ 6 \\ 7 \\ 3150 \\ 100 \end{pmatrix}
\]

$\vx_0 + (1,\, 2,\, 3)^T$ : \textbf{non}, car $\vx_0 \in \mathbb{R}^5$ et $(1,2,3)^T \in \mathbb{R}^3$.
L'addition n'est d\'{e}finie que pour des vecteurs de \textbf{m\^{e}me dimension}.

\begin{attentionbox}[Erreurs fr\'{e}quentes]
    \begin{itemize}
        \item Confondre la \textbf{dimension} du vecteur ($p = 5$) avec le \textbf{nombre d'observations} ($n = 30$).
        \item Oublier la \textbf{transpos\'{e}e} : un vecteur colonne $\vx$ et sa transpos\'{e}e $\vx^T$ n'ont pas le m\^{e}me shape en NumPy.
    \end{itemize}
\end{attentionbox}

% ======================================================================
% EXERCICE 2
% ======================================================================
\section*{Exercice 2 \hfill \exofondamental}
\addcontentsline{toc}{section}{Exercice 2 --- Produit scalaire \`{a} la main}
\textit{Rappel : $\pscal{\vx}{\vw} = \sum_{i=1}^{p} x_i w_i$.}

\separateur

\textbf{a) Calcul de $\pscal{\vx}{\vw}$ :}
\begin{align*}
    \pscal{\vx}{\vw} &= 3 \times 0{,}5 + 5 \times 0{,}3 + 2 \times 0{,}8 \\
                     &= 1{,}5 + 1{,}5 + 1{,}6 \\
                     &= \boxed{4{,}6}
\end{align*}

\textbf{b) Pr\'{e}diction pour le d\'{e}veloppeur \#0} avec
$\vw = (2{,}0;\; 1{,}5;\; 3{,}0;\; 0{,}005;\; 0{,}02)^T$ et $b = 15$ :
\begin{align*}
    \pscal{\vx_0}{\vw} &= 0 \times 2{,}0 + 2 \times 1{,}5 + 3 \times 3{,}0 + 150 \times 0{,}005 + 80 \times 0{,}02 \\
                       &= 0 + 3{,}0 + 9{,}0 + 0{,}75 + 1{,}6 = 14{,}35
\end{align*}
Pr\'{e}diction : $\yhat_0 = 14{,}35 + 15 = \boxed{29{,}35}$ k\texteuro.
Salaire r\'{e}el : $y_0 = 28{,}0$ k\texteuro. Erreur : $|29{,}35 - 28{,}0| = 1{,}35$ k\texteuro.

\textbf{c) Pi\`{e}ge --- dimensions incompatibles.}

\textbf{Non}, on ne peut pas calculer ce produit scalaire.
Le premier vecteur est dans $\mathbb{R}^5$ et le second dans $\mathbb{R}^3$.
Le produit scalaire exige que les deux vecteurs aient la \textbf{m\^{e}me dimension}.

\begin{attentionbox}[Erreurs fr\'{e}quentes]
    \begin{itemize}
        \item Oublier d'additionner le \textbf{biais}~$b$ apr\`{e}s le produit scalaire.
        \item Erreur de calcul sur les termes avec des petits poids ($0{,}005 \times 150$ se calcule mieux comme $150 / 200 = 0{,}75$).
        \item Tenter de calculer un produit scalaire avec des dimensions diff\'{e}rentes --- en NumPy, cela provoque une \texttt{ValueError}.
    \end{itemize}
\end{attentionbox}

% ======================================================================
% EXERCICE 3
% ======================================================================
\section*{Exercice 3 \hfill \exofondamental}
\addcontentsline{toc}{section}{Exercice 3 --- Multiplication matrice-vecteur}
\textit{Rappel : $(\mX\vw)_i = \vx_i^T \vw$ (produit scalaire de la ligne $i$ avec $\vw$).}

\separateur

\textbf{a) Dimensions.}
$\mX \in \mathbb{R}^{3 \times 2}$, $\vw \in \mathbb{R}^{2}$,
donc $\hat{\vy} = \mX\vw \in \mathbb{R}^{3}$ (3 pr\'{e}dictions).

\textbf{b) Calcul ligne par ligne :}
\begin{align*}
    \yhat_1 &= 1 \times 2 + 3 \times 0{,}5 = 2 + 1{,}5 = \boxed{3{,}5} \\
    \yhat_2 &= 4 \times 2 + 2 \times 0{,}5 = 8 + 1{,}0 = \boxed{9{,}0} \\
    \yhat_3 &= 0 \times 2 + 5 \times 0{,}5 = 0 + 2{,}5 = \boxed{2{,}5}
\end{align*}
\[
    \hat{\vy} = \mX\vw = \begin{pmatrix} 3{,}5 \\ 9{,}0 \\ 2{,}5 \end{pmatrix}
\]

\textbf{c) V\'{e}rification :}
$\vx_1^T = (1 \;\; 3)$, donc $\vx_1^T \vw = 1 \times 2 + 3 \times 0{,}5 = 3{,}5 = \yhat_1$. \checkmark

\begin{retenirbox}
    La multiplication $\mX\vw$ revient \`{a} $n$ produits scalaires en parall\`{e}le.
    Chaque ligne de $\mX$ fait un produit scalaire avec $\vw$.
\end{retenirbox}

% ======================================================================
% EXERCICE 4
% ======================================================================
\section*{Exercice 4 \hfill \exofondamental}
\addcontentsline{toc}{section}{Exercice 4 --- Calcul du MSE}
\textit{Rappel : $\text{MSE} = \frac{1}{n}\sum_{i=1}^{n}(y_i - \yhat_i)^2$.}

\separateur

\textbf{a) R\'{e}sidus $e_i = y_i - \yhat_i$ :}

\begin{center}
\renewcommand{\arraystretch}{1.2}
\begin{tabular}{c|c|c|c|c}
    \toprule
    $i$ & $y_i$ & $\yhat_i$ & $e_i$ & $e_i^2$ \\
    \midrule
    1 & 28.0 & 31.3 & $-3{,}3$  & $10{,}89$ \\
    2 & 35.0 & 33.3 & $+1{,}7$  & $2{,}89$  \\
    3 & 38.0 & 35.4 & $+2{,}6$  & $6{,}76$  \\
    4 & 44.0 & 41.6 & $+2{,}4$  & $5{,}76$  \\
    5 & 52.0 & 49.9 & $+2{,}1$  & $4{,}41$  \\
    \midrule
    & & & \textbf{Somme} & $\mathbf{30{,}71}$ \\
    \bottomrule
\end{tabular}
\end{center}

\textbf{b-c) Calcul du MSE :}
\formule{\text{MSE} = \frac{1}{5} \times 30{,}71 = \boxed{6{,}14 \;\text{(k\texteuro)}^2}}

\textbf{d) Erreur typique :}
$\sqrt{\text{MSE}} = \sqrt{6{,}14} \approx \boxed{2{,}48 \text{ k\texteuro}}$.
Le mod\`{e}le se trompe en moyenne d'environ 2\,500\,\texteuro{} --- c'est raisonnable.

\begin{attentionbox}[Erreurs fr\'{e}quentes]
    \begin{itemize}
        \item Oublier de mettre au \textbf{carr\'{e}} avant de sommer (sinon les erreurs $+$ et $-$ s'annulent).
        \item Diviser par $n-1$ au lieu de $n$ (c'est la variance \textit{non biais\'{e}e}, pas le MSE).
        \item Le MSE est en \textbf{(k\texteuro)$^2$}, pas en k\texteuro{}. C'est $\sqrt{\text{MSE}}$ qui est en k\texteuro.
    \end{itemize}
\end{attentionbox}

% ======================================================================
% EXERCICE 5
% ======================================================================
\section*{Exercice 5 \hfill \exointermediaire}
\addcontentsline{toc}{section}{Exercice 5 --- Interpr\'{e}ter les poids}
\textit{Mod\`{e}le : $\yhat = 25 + 3 \times \texttt{exp} + 1{,}5 \times \texttt{lang}$.}

\separateur

\textbf{a) Identification des param\`{e}tres :}
\begin{itemize}
    \item $w_0 = 25$ (biais) : salaire de base pr\'{e}dit quand exp\'{e}rience $= 0$ et langages $= 0$.
    \item $w_1 = 3$ : chaque ann\'{e}e d'exp\'{e}rience suppl\'{e}mentaire augmente le salaire pr\'{e}dit de \textbf{3\,k\texteuro}.
    \item $w_2 = 1{,}5$ : chaque langage suppl\'{e}mentaire augmente le salaire pr\'{e}dit de \textbf{1,5\,k\texteuro}.
\end{itemize}

\textbf{b)} $\yhat = 25 + 3 \times 0 + 1{,}5 \times 1 = \boxed{26{,}5}$ k\texteuro.
C'est r\'{e}aliste pour un d\'{e}butant avec un seul langage en France.

\textbf{c)} $+1$ an d'exp\'{e}rience $\Rightarrow$ $+3$\,k\texteuro.
$+1$ langage $\Rightarrow$ $+1{,}5$\,k\texteuro.
L'exp\'{e}rience a \textbf{deux fois plus d'impact} que le nombre de langages.

\textbf{d)} $\yhat = 25 + 3 \times 5 + 1{,}5 \times 4 = 25 + 15 + 6 = \boxed{46}$ k\texteuro.
Le d\'{e}veloppeur \#9 a un vrai salaire de 44\,k\texteuro{} --- erreur de 2\,k\texteuro, le mod\`{e}le est assez bon ici.

\begin{retenirbox}
    Les poids d'un mod\`{e}le lin\'{e}aire s'interpr\`{e}tent comme l'\textbf{impact marginal}
    de chaque feature sur la pr\'{e}diction, \emph{toutes choses \'{e}gales par ailleurs}.
\end{retenirbox}

% ======================================================================
% EXERCICE 6
% ======================================================================
\section*{Exercice 6 \hfill \exointermediaire}
\addcontentsline{toc}{section}{Exercice 6 --- Choisir le meilleur mod\`{e}le}
\textit{Mod\`{e}le A : $\yhat = 2x + 32$. Mod\`{e}le B : $\yhat = 3x + 25$.}

\separateur

\textbf{a-b) Pr\'{e}dictions :}

\begin{center}
\renewcommand{\arraystretch}{1.2}
\begin{tabular}{c|c|c|c|c|c}
    \toprule
    \# & exp & $y_i$ & $\yhat_i^A = 2x{+}32$ & $\yhat_i^B = 3x{+}25$ & $(e_i^A)^2$ \;\; $(e_i^B)^2$ \\
    \midrule
    0 & 0 & 28.0 & 32.0 & 25.0 & 16.00 \;\;\; 9.00 \\
    1 & 1 & 26.5 & 34.0 & 28.0 & 56.25 \;\;\; 2.25 \\
    2 & 1 & 35.0 & 34.0 & 28.0 &  1.00 \;\;\; 49.00 \\
    3 & 2 & 32.0 & 36.0 & 31.0 & 16.00 \;\;\; 1.00 \\
    4 & 2 & 38.0 & 36.0 & 31.0 &  4.00 \;\;\; 49.00 \\
    5 & 0 & 25.0 & 32.0 & 25.0 & 49.00 \;\;\; 0.00 \\
    \bottomrule
\end{tabular}
\end{center}

\textbf{c) MSE :}
\begin{align*}
    \text{MSE}_A &= \frac{16{,}00 + 56{,}25 + 1{,}00 + 16{,}00 + 4{,}00 + 49{,}00}{6}
    = \frac{142{,}25}{6} = \boxed{23{,}71} \\[4pt]
    \text{MSE}_B &= \frac{9{,}00 + 2{,}25 + 49{,}00 + 1{,}00 + 49{,}00 + 0{,}00}{6}
    = \frac{110{,}25}{6} = \boxed{18{,}38}
\end{align*}

\textbf{d)} Le mod\`{e}le B est meilleur (MSE plus faible). Cependant, ni A ni B n'est optimal :
la droite des moindres carr\'{e}s ($\yhat = 2{,}07x + 31{,}27$) a un MSE de 20{,}46 sur les 30 observations
compl\`{e}tes. Sur un si petit \'{e}chantillon (6 observations), le classement peut changer.

\begin{attentionbox}[Erreurs fr\'{e}quentes]
    \begin{itemize}
        \item Comparer les MSE de mod\`{e}les \'{e}valu\'{e}s sur des \'{e}chantillons \textbf{diff\'{e}rents} --- il faut toujours utiliser les m\^{e}mes donn\'{e}es.
        \item Conclure qu'un mod\`{e}le est ``le meilleur'' alors qu'il n'est meilleur que sur ce sous-ensemble.
    \end{itemize}
\end{attentionbox}

% ======================================================================
% EXERCICE 7
% ======================================================================
\section*{Exercice 7 \hfill \exointermediaire}
\addcontentsline{toc}{section}{Exercice 7 --- Matrice et dataset}
\textit{Construire $\mX$ et $\vy$ \`{a} partir de profils textuels.}

\separateur

\textbf{a) Matrice de features $\mX \in \mathbb{R}^{4 \times 5}$ :}

\emph{Encodage :} Bac$=1$, Bac+3$=3$, Bac+5$=4$.

\formule{
    \mX = \begin{pmatrix}
        3  & 2 & 3 & 150  & 80 \\
        7  & 5 & 4 & 800  & 50 \\
        0  & 1 & 1 & 20   & 100 \\
        12 & 6 & 4 & 3000 & 30
    \end{pmatrix}
    \in \mathbb{R}^{4 \times 5}
}

\textbf{b) Vecteur cible :}
$\vy = \begin{pmatrix} 34 & 50 & 24 & 58 \end{pmatrix}^T \in \mathbb{R}^4$.

\textbf{c)} $x_{2,4} = 100$. C'est le pourcentage de t\'{e}l\'{e}travail de Clara
(2\textsuperscript{e} ligne, 4\textsuperscript{e} feature en partant de 0 --- ou
3\textsuperscript{e} ligne, 5\textsuperscript{e} colonne en indexation \`{a} partir de 1).

\begin{attentionbox}[Attention \`{a} l'indexation]
    L'indexation math\'{e}matique commence \`{a} \textbf{1} ($x_{1,1}$) tandis que Python/NumPy
    commence \`{a} \textbf{0} (\texttt{X[0,0]}). Ici Bob est la ligne 2 en math (ou \texttt{X[1]} en Python).
\end{attentionbox}

\textbf{d) Pr\'{e}diction pour Bob :}
$\vx_{\text{Bob}} = (7,\; 5,\; 4,\; 800,\; 50)^T$.
\begin{align*}
    \vw^T \vx_{\text{Bob}} &= 7 \times 2 + 5 \times 1{,}5 + 4 \times 3 + 800 \times 0{,}005 + 50 \times 0{,}02 \\
    &= 14 + 7{,}5 + 12 + 4 + 1 = 38{,}5
\end{align*}
$\yhat_{\text{Bob}} = 38{,}5 + 15 = \boxed{53{,}5}$ k\texteuro. Vrai salaire : 50\,k\texteuro{} --- erreur de 3{,}5\,k\texteuro.

% ======================================================================
% EXERCICE 8
% ======================================================================
\section*{Exercice 8 \hfill \exointermediaire}
\addcontentsline{toc}{section}{Exercice 8 --- Vecteurs NumPy (Python guid\'{e})}
\textit{Cr\'{e}er des vecteurs, calculer norme et produit scalaire.}

\separateur

\begin{pythonbox}[Correction compl\`{e}te]
\begin{lstlisting}
import numpy as np

# 1. Vecteur du dev #3 : exp=2, lang=2, etudes=3, entrep=200, remote=60
x3 = np.array([2, 2, 3, 200, 60])

# 2. Dimension
print(f"Dimension de x3 : {x3.shape}")  # (5,)

# 3. 3e composante (index 2 en Python)
niveau = x3[2]
print(f"Niveau d'etudes : {niveau}")  # 3

# 4. Norme euclidienne
norme = np.sqrt(np.sum(x3 ** 2))
print(f"Norme de x3 : {norme:.2f}")  # 209.11

# 5. Vecteur du dev #5 : exp=0, lang=1, etudes=1, entrep=30, remote=50
x5 = np.array([0, 1, 1, 30, 50])

# 6. Produit scalaire avec @
ps = x3 @ x5
print(f"Produit scalaire x3 . x5 = {ps}")  # 9185

# 7. Distance euclidienne
distance = np.sqrt(np.sum((x3 - x5) ** 2))
print(f"Distance : {distance:.2f}")  # 170.27
\end{lstlisting}
\end{pythonbox}

\textbf{Sortie attendue :}
\begin{verbatim}
Dimension de x3 : (5,)
Niveau d'etudes : 3
Norme de x3 : 209.11
Produit scalaire x3 . x5 = 9185
Distance : 170.27
\end{verbatim}

\begin{intuitionbox}[Variante]
    La norme peut aussi se calculer avec \texttt{np.linalg.norm(x3)}.
    Le produit scalaire avec \texttt{np.dot(x3, x5)}.
    L'op\'{e}rateur \texttt{@} est le plus lisible et le plus courant en pratique.
\end{intuitionbox}

\begin{attentionbox}[Erreurs fr\'{e}quentes en Python]
    \begin{itemize}
        \item \texttt{x3[3]} donne la 4\textsuperscript{e} composante (index 3), pas la 3\textsuperscript{e} --- l'indexation commence \`{a}~0.
        \item La norme de $\vx$ est domin\'{e}e par les grandes valeurs (ici 200 et 60). C'est pourquoi on \textbf{normalise} souvent les features.
    \end{itemize}
\end{attentionbox}

% ======================================================================
% EXERCICE 9
% ======================================================================
\section*{Exercice 9 \hfill \exointermediaire}
\addcontentsline{toc}{section}{Exercice 9 --- Matrice du dataset (Python guid\'{e})}
\textit{Charger les donn\'{e}es dans une matrice, slicing lignes/colonnes.}

\separateur

\begin{pythonbox}[Correction compl\`{e}te]
\begin{lstlisting}
import numpy as np

# 1. Matrice des 6 premiers developpeurs (5 features)
X = np.array([
    [0, 2, 3, 150, 80],     # dev #0
    [1, 1, 2, 50, 100],     # dev #1
    [1, 3, 4, 800, 40],     # dev #2
    [2, 2, 3, 200, 60],     # dev #3
    [2, 4, 4, 3000, 20],    # dev #4
    [0, 1, 1, 30, 50],      # dev #5
])

# 2. Verifier les dimensions
print(f"Shape de X : {X.shape}")     # (6, 5)

# 3. Colonne "experience" = colonne d'index 0
experience = X[:, 0]
print(f"Experiences : {experience}")  # [0 1 1 2 2 0]

# 4. Ligne du developpeur #4 = index 4
dev4 = X[4, :]
print(f"Developpeur #4 : {dev4}")    # [2 4 4 3000 20]

# 5. Vecteur des salaires
y = np.array([28.0, 26.5, 35.0, 32.0, 38.0, 25.0])

# 6. Transposee
print(f"Shape de X.T : {X.T.shape}") # (5, 6)

# 7. Element X[2, 3] = dev #2, feature #3 (taille entreprise)
print(f"X[2,3] = {X[2, 3]}")         # 800
# C'est la taille de l'entreprise du developpeur #2
\end{lstlisting}
\end{pythonbox}

\textbf{Sortie attendue :}
\begin{verbatim}
Shape de X : (6, 5)
Experiences : [0 1 1 2 2 0]
Developpeur #4 : [   2    4    4 3000   20]
Shape de X.T : (5, 6)
X[2,3] = 800
\end{verbatim}

\begin{attentionbox}[Erreurs fr\'{e}quentes en Python]
    \begin{itemize}
        \item Confondre \texttt{X[2]} (toute la ligne 2) et \texttt{X[:, 2]} (toute la colonne 2).
        \item Oublier le \texttt{:} dans le slicing : \texttt{X[4]} et \texttt{X[4, :]} sont \'{e}quivalents, mais \texttt{X[:, 0]} est la seule fa\c{c}on d'extraire une colonne.
    \end{itemize}
\end{attentionbox}

% ======================================================================
% EXERCICE 10
% ======================================================================
\section*{Exercice 10 \hfill \exointermediaire}
\addcontentsline{toc}{section}{Exercice 10 --- R\'{e}gression from scratch (Python guid\'{e})}
\textit{Calculer la droite des moindres carr\'{e}s, le MSE, et tracer le graphique.}

\separateur

\begin{pythonbox}[Correction compl\`{e}te]
\begin{lstlisting}
import numpy as np
import matplotlib.pyplot as plt

experience = np.array([0,1,1,2,2,0,3,1,4,5,5,6,6,7,7,4,
                        8,5,9,10,10,11,12,9,13,11,15,17,20,18])
salaire = np.array([28.0,26.5,35.0,32.0,38.0,25.0,34.5,42.0,
    38.5,44.0,41.0,48.0,43.5,50.0,39.0,47.0,51.0,40.0,
    52.0,55.0,48.5,58.0,60.0,42.0,62.0,54.0,65.0,68.0,
    72.0,58.0])

# 1. Moyennes
x_mean = experience.mean()          # 7.40
y_mean = salaire.mean()             # 46.58

# 2. Pente w (formule des moindres carres)
numerateur = np.sum((experience - x_mean) * (salaire - y_mean))
denominateur = np.sum((experience - x_mean) ** 2)
w = numerateur / denominateur       # 2.07
print(f"Pente w = {w:.4f}")

# 3. Biais b
b = y_mean - w * x_mean             # 31.27
print(f"Biais b = {b:.4f}")

# 4. Predictions
y_pred = w * experience + b

# 5. MSE
mse = np.mean((salaire - y_pred) ** 2)  # 20.46
print(f"MSE = {mse:.2f}")

# 6. Graphique
fig, ax = plt.subplots(figsize=(8, 5))
ax.scatter(experience, salaire, c='blue', label='Donnees')
x_line = np.linspace(0, 21, 100)
ax.plot(x_line, w * x_line + b, 'r--',
        label=f'y={w:.2f}x+{b:.2f}')
ax.set_xlabel("Experience (annees)")
ax.set_ylabel("Salaire (kEUR)")
ax.set_title("Regression lineaire from scratch")
ax.legend(); ax.grid(True, alpha=0.3)
plt.show()
\end{lstlisting}
\end{pythonbox}

\textbf{Sortie attendue :}
\begin{verbatim}
Pente w = 2.0672
Biais b = 31.2830
MSE = 20.46
\end{verbatim}
+ un graphique montrant le nuage de points bleu et la droite rouge en pointill\'{e}s.

\begin{intuitionbox}[Variante : avec scikit-learn]
    En une ligne : \texttt{from sklearn.linear\_model import LinearRegression}.
    Mais l'objectif ici est de comprendre le m\'{e}canisme \textbf{avant} d'utiliser la biblioth\`{e}que.
\end{intuitionbox}

\begin{retenirbox}
    La formule de la pente $w_1 = \frac{\sum(x_i - \bar{x})(y_i - \bar{y})}{\sum(x_i - \bar{x})^2}$
    mesure \`{a} quel point $x$ et $y$ varient \textbf{ensemble}, normalis\'{e} par la variance de~$x$.
\end{retenirbox}

% ======================================================================
% EXERCICE 11
% ======================================================================
\section*{Exercice 11 \hfill \exoavance}
\addcontentsline{toc}{section}{Exercice 11 --- Explorer le dataset (Python autonome)}
\textit{Statistiques descriptives, corr\'{e}lations et visualisation.}

\separateur

\begin{pythonbox}[Correction compl\`{e}te]
\begin{lstlisting}
import numpy as np
import matplotlib.pyplot as plt

# --- Chargement du dataset ---
data = np.array([
    [ 0,2,3, 150, 80,28.0],[ 1,1,2,  50,100,26.5],
    [ 1,3,4, 800, 40,35.0],[ 2,2,3, 200, 60,32.0],
    [ 2,4,4,3000, 20,38.0],[ 0,1,1,  30, 50,25.0],
    [ 3,3,3, 100, 80,34.5],[ 1,5,4,  15,100,42.0],
    [ 4,3,3, 500, 60,38.5],[ 5,4,4,1200, 40,44.0],
    [ 5,3,3, 300, 80,41.0],[ 6,5,4,2000, 30,48.0],
    [ 6,4,3, 150, 90,43.5],[ 7,5,4, 800, 50,50.0],
    [ 7,3,2,  80, 70,39.0],[ 4,6,5, 400,100,47.0],
    [ 8,4,4,1500, 40,51.0],[ 5,2,3,4000, 10,40.0],
    [ 9,5,4, 600, 60,52.0],[10,6,4,2500, 30,55.0],
    [10,4,3, 200, 80,48.5],[11,5,5,1000, 50,58.0],
    [12,7,4,3500, 20,60.0],[ 9,3,3,  60, 90,42.0],
    [13,6,4, 800, 70,62.0],[11,4,4,5000, 10,54.0],
    [15,8,4,1200, 50,65.0],[17,6,5,2000, 40,68.0],
    [20,7,4, 500, 60,72.0],[18,5,3, 100, 80,58.0],
])

# a) Dimensions et apercu
print(f"Dimensions : {data.shape}")
print("5 premieres lignes :")
print(data[:5])

# b) Statistiques par colonne
noms = ["experience","nb_langages","niveau_etudes",
        "taille_entreprise","remote","salaire"]
print("\n--- Statistiques descriptives ---")
for j, nom in enumerate(noms):
    col = data[:, j]
    print(f"{nom:20s} | moy={col.mean():8.2f} "
          f"std={col.std():8.2f} "
          f"min={col.min():7.1f} max={col.max():7.1f}")

# c) Correlations avec le salaire
salaire = data[:, 5]
print("\n--- Correlations avec le salaire ---")
for j in range(5):
    r = np.corrcoef(data[:, j], salaire)[0, 1]
    print(f"{noms[j]:20s} | r = {r:.3f}")

# d) 5 sous-figures
fig, axes = plt.subplots(1, 5, figsize=(18, 3.5))
for j in range(5):
    r = np.corrcoef(data[:, j], salaire)[0, 1]
    axes[j].scatter(data[:, j], salaire, s=30, alpha=0.7)
    axes[j].set_xlabel(noms[j])
    axes[j].set_ylabel("salaire" if j == 0 else "")
    axes[j].set_title(f"r = {r:.2f}")
    axes[j].grid(True, alpha=0.3)
plt.suptitle("Correlations features vs salaire")
plt.tight_layout(); plt.show()
\end{lstlisting}
\end{pythonbox}

\textbf{R\'{e}sultats cl\'{e}s :}
\begin{center}
\begin{tabular}{l c l}
    \toprule
    \textbf{Feature} & $r$ & \textbf{Interpr\'{e}tation} \\
    \midrule
    experience        & $0{,}927$ & tr\`{e}s forte corr\'{e}lation positive \\
    nb\_langages      & $0{,}897$ & forte corr\'{e}lation positive \\
    niveau\_etudes    & $0{,}667$ & corr\'{e}lation mod\'{e}r\'{e}e \\
    taille\_entreprise & $0{,}264$ & faible corr\'{e}lation \\
    remote            & $-0{,}341$ & faible corr\'{e}lation n\'{e}gative \\
    \bottomrule
\end{tabular}
\end{center}

L'exp\'{e}rience est la feature la plus corr\'{e}l\'{e}e au salaire, suivie du nombre de langages.

\begin{attentionbox}[Erreurs fr\'{e}quentes]
    \begin{itemize}
        \item Utiliser \texttt{np.corrcoef} sans extraire l'\'{e}l\'{e}ment \texttt{[0,1]} --- la fonction renvoie une \textbf{matrice} $2 \times 2$, pas un scalaire.
        \item Confondre corr\'{e}lation et causalit\'{e} (le remote n\'{e}gatif ne signifie pas que le t\'{e}l\'{e}travail \textit{r\'{e}duit} le salaire).
    \end{itemize}
\end{attentionbox}

% ======================================================================
% EXERCICE 12
% ======================================================================
\section*{Exercice 12 \hfill \exoavance}
\addcontentsline{toc}{section}{Exercice 12 --- Tester plusieurs mod\`{e}les (Python autonome)}
\textit{Comparer 3 droites de r\'{e}gression, calculer le MSE, tracer le graphique.}

\separateur

\begin{pythonbox}[Correction compl\`{e}te]
\begin{lstlisting}
import numpy as np
import matplotlib.pyplot as plt

# --- Donnees ---
experience = np.array([0,1,1,2,2,0,3,1,4,5,5,6,6,7,7,4,
    8,5,9,10,10,11,12,9,13,11,15,17,20,18])
salaire = np.array([28.0,26.5,35.0,32.0,38.0,25.0,34.5,
    42.0,38.5,44.0,41.0,48.0,43.5,50.0,39.0,47.0,51.0,
    40.0,52.0,55.0,48.5,58.0,60.0,42.0,62.0,54.0,65.0,
    68.0,72.0,58.0])

# --- 3 modeles ---
modeles = {
    "A": (1.5, 33),
    "B": (2.07, 31.27),
    "C": (2.8, 28),
}

# a-b) Predictions et MSE
print("Modele |  w1   |  w0   |   MSE")
print("-------|-------|-------|--------")
resultats = {}
for nom, (w1, w0) in modeles.items():
    y_pred = w1 * experience + w0
    mse = np.mean((salaire - y_pred) ** 2)
    resultats[nom] = (w1, w0, mse)
    print(f"   {nom}   | {w1:5.2f} | {w0:5.2f} | {mse:.2f}")

# c) Graphique
fig, ax = plt.subplots(figsize=(9, 5))
ax.scatter(experience, salaire, c='#333', s=40,
           zorder=3, label='Donnees')

couleurs = {"A": "#e74c3c", "B": "#2ecc71", "C": "#3498db"}
x_line = np.linspace(0, 21, 100)
for nom, (w1, w0, mse) in resultats.items():
    ax.plot(x_line, w1 * x_line + w0,
            color=couleurs[nom], lw=2,
            label=f'{nom}: y={w1}x+{w0} (MSE={mse:.1f})')

ax.set_xlabel("Experience (annees)")
ax.set_ylabel("Salaire (kEUR)")
ax.set_title("Comparaison de 3 modeles lineaires")
ax.legend(); ax.grid(True, alpha=0.3)
plt.tight_layout(); plt.show()

# d) Conclusion
print("\nLe modele B a le MSE le plus faible.")
print("C'est normal : B est la droite des moindres")
print("carres, qui minimise le MSE par construction.")
\end{lstlisting}
\end{pythonbox}

\textbf{Sortie attendue :}
\begin{verbatim}
Modele |  w1   |  w0   |   MSE
-------|-------|-------|--------
   A   |  1.50 | 33.00 | 27.03
   B   |  2.07 | 31.27 | 20.46
   C   |  2.80 | 28.00 | 21.31
\end{verbatim}

Le mod\`{e}le \textbf{B} a le MSE le plus faible ($20{,}46$), car c'est pr\'{e}cis\'{e}ment la droite
des moindres carr\'{e}s --- celle qui, \textbf{par d\'{e}finition}, minimise le MSE.

Le mod\`{e}le C a un MSE tr\`{e}s proche ($21{,}31$) : une pente un peu trop forte compens\'{e}e par un biais plus bas.
Le mod\`{e}le A ($27{,}03$) sous-estime l'impact de l'exp\'{e}rience (pente trop faible).

\begin{retenirbox}[Conclusion g\'{e}n\'{e}rale]
    La droite des moindres carr\'{e}s est \textbf{optimale} par construction : aucun autre jeu de poids
    $(w_1, w_0)$ ne peut donner un MSE plus faible sur les m\^{e}mes donn\'{e}es. C'est ce
    que nous d\'{e}montrerons formellement au \textbf{Jour 2} en introduisant l'\textbf{optimisation}
    et la \textbf{descente de gradient}.
\end{retenirbox}

\begin{intuitionbox}[Variante : boucle automatique]
    On peut automatiser la recherche en testant des milliers de poids :
    c'est l'id\'{e}e de la \textbf{descente de gradient} du Jour 2.
\end{intuitionbox}

% ======================================================================
\vfill
\begin{center}
    {\footnotesize\sffamily\textcolor{gray}{%
        \textcopyright\ 2025 Mohamed EL AFRIT --- IPSSI \enspace\textbullet\enspace
        \url{www.mohamedelafrit.com} \enspace\textbullet\enspace
        CC BY-NC-SA 4.0
    }}
\end{center}

\end{document}
